%%%%%%%%%%%%%%%%%%%%%%%%%%%%%%%%%%%%%%%%%%%%%%%%%%%%%%%%%
% Niniejszy plik przedstawia przykładowy skład pracy dyplomowej na Wydziale Matematyki PWr. 
% 
% Autorzy: Damian Fafuła, Michał Kijaczko, Jakub Michalczak, Maciej Miśta Dagmara Nowak, Tomasz Skalski, Wojciech Słomian
% Data utworzenia: 8.05.2018
%
% Zmodyfikowane przez: Jakuba Ślęzaka 
% Data utworzenia aktualnej wersji: 28.01.2025 
% Numer wersji: 1.1
%
% Poniższą formatkę można rozpowszechniać i edytować pod warunkiem zachowania numeru wersji, 
% informacji o autorach i dodaniu informacji o wprowadzonych zmianach.
%
%                          OPCJE
%%%%%%%%%%%%%%%%%%%%%%%%%%%%%%%%%%%%%%%%%%%%%%%%%%%%%%%%%

% Następna linia zaczyna kolejne rozdziały od stron nieparzystych
\def\rodzialyOdNieparzystych {wybrane}
% zakomentować ją, jeżeli rozdziały są krótkie i nie jest to wskazane.

% Następna linia zaznacza kolorami interaktywne linki w pliku pdf
\def\koloroweLinki {wybrane}
% zakomentować ją, aby były czarno białe, np. do druku

% Opcje strony tytułowej:
% Domyślną opcją jest: praca magisterska, język polski.
% W przypadku pracy pisanej w języku angielskim dodajemy opcję [english].
% Dla pracy licencjackiej dodajemy opcję [licencjacka].
% Dla pracy inżynierskiej dodajemy opcję [inzynierska].
% Dopuszczalne są podwójne opcje, np. [licencjacka, english].
% Opcje dodajemy w kwadratowym nawiasie przy \documentclass.
%
%
%%%%%%%%%%%%%%%%%%%%%%%%%%%%%%%%%%%%%%%%%%%%%%%%%%%%%%%%%
\documentclass[]{pwr_wmat_praca_dyplomowa}


%                       DANE DO PRACY
%%%%%%%%%%%%%%%%%%%%%%%%%%%%%%%%%%%%%%%%%%%%%%%%%%%%%%%%%

% W przypadku pracy dyplomowej w języku angielskim nie jest konieczne 
% wypełnianie pól: \tytul{}, \kierunek{}, \specjalnosc{}, \streszczenie{}, \slowakluczowe{}.
%
% Imię i nazwisko autora
\autor{Imię i nazwisko dyplomanta}
%
% Tytuł pracy dyplomowej 
\tytul{Tytuł pracy dyplomowej} 
\tytulang{Tytuł pracy dyplomowej w języku angielskim}
%
% Tytuł / stopień / imię i nazwisko opiekuna
\opiekun{prof. dr hab. Jan Kowalski}

% Kierunek studiów
\kierunekstudiow{Matematyka}

% Kierunek studiów po angielsku
\kierunekstudiowang{Mathematics}

% Specjalność
\specjalnosc{Matematyka teoretyczna} 

% Specjalność w języku angielskim
\specjalnoscang{Theoretical Mathematics} 

% Krótkie streszczenia po polsku i angielsku, nie dłuższe niż 530 znaków.
\streszczenie{Tutaj piszemy krótkie streszczenie pracy (nie powinno być dłuższe niż 530 znaków).}
\streszczenieang{Tutaj piszemy krótkie streszczenie pracy w języku angielskim (nie powinno być dłuższe niż 530 znaków).}
%
% Podajemy najważniejsze słowa kluczowe po polsku i angielsku, w obu przypadkach, max 150 znaków.
\slowakluczowe{tutaj podajemy najważniejsze słowa kluczowe (łącznie nie powinny być dłuższe niż 150 znaków).}  
\slowakluczoweang{tutaj podajemy najważniejsze słowa kluczowe w języku angielskim (łącznie nie powinny być dłuższe niż 150 znaków)}
%
%
%%%%%%%%%%%%%%%%%%%%%%%%%%%%%%%%%%%%%%%%%%%%%%%%%%%%%%%%%

% Definicje komend tworzących twierdzenia, definicje, itd, można zmienić wedle uznania.
\theoremstyle{plain}
\newtheorem{theorem}{Twierdzenie}
\numberwithin{theorem}{chapter}
\newtheorem{lemma}[theorem]{Lemat} 
\newtheorem{corollary}[theorem]{Wniosek}
\newtheorem{fact}[theorem]{Fakt}
\theoremstyle{definition}
\numberwithin{theorem}{chapter}
\newtheorem{definition}[theorem]{Definicja} 
\newtheorem{example}[theorem]{Przykład}
\newtheorem{note}[theorem]{Uwaga}



%                     POCZĄTEK PRACY
%%%%%%%%%%%%%%%%%%%%%%%%%%%%%%%%%%%%%%%%%%%%%%%%%%%%%%%%%%%%%%%%%
\begin{document}
\frontmatter
\maketitle
\mainmatter
\tableofcontents
%\listoffigures
%\listoftables

\chapter{Wstęp}
We wstępie zapowiadamy, o czym będzie praca. Próbujemy zachęcić czytelnika do dalszej lektury, wprowadzamy w dziecinę pracy, wyjaśniamy, dlaczego wybraliśmy właśnie ten temat i co nas w nim zainteresowało.

\chapter{Główna część pracy}

\section{Przykładowa tabela}
Tabela \ref{tab:przykladowa} przedstawia przykładową tabelę. Do tworzenia tabeli służą m.in. środowiska \texttt{tabular} oraz \texttt{table}. Tytuł powinien znajdować się centralnie nad tabelą. Z opisu tabeli powinno być jasne źródło przedstawianych danych np. ,,opracowano na podstawie danych z GUS``, ,,wyniki symulacji przeprowadzonych metodą X'', itp.
\begin{table}[ht]
\caption{Podstawowa Tabela}
\centering
\begin{tabular}{ccc}
	\hline
	\hline                       
	Państwo & PKB (w milionach USD )& Stopa bezrobocia  \\  [0.5ex] 
	\hline 
	Stany Zjednoczone & 75 278 049 & 4,60\%  \\
	Chiny & 11 218 281 & 4,10\%   \\
	Japonia & 4 938 644 & 3,10\%  \\
	Niemcy & 3 466 639 & 6,00\%   \\
	Wielka Brytania & 2 629 188 & 4,60\%  \\ [1ex]  
	\hline 
\end{tabular}
\caption*{\textit{Źródło danych: GUS}} \label{tab:przykladowa} 
\end{table}




\section{Przykładowy rysunek}

Rysunki do pracy dyplomowej należy wstawiać w sposób podobny do wstawiania tabel, z~zasadniczą różnicą polegającą na tym, że podpis powinno umieszczać się centralnie pod rysunkiem, a nie powyżej niego. Numeracja i sposób cytowania pozostają bez zmian, przy czym tabele i rysunki nie mają numeracji wspólnej, np. po Tabeli \ref{tab:przykladowa} występuje Rysunek \ref{fig:WMat} (o ile jest to pierwszy rysunek rozdziału pierwszego), a nie Rysunek $1.3$.

\begin{figure}[!ht] \centering
                     
\includegraphics[scale=0.27]{fig/logo_w13.jpg}

\caption{Logo WMat.} \label{fig:WMat} % etykietę umieszczamy zawsze za \caption
\end{figure}


\section{Przykładowe równania}

Przykładowe równanie jednoliniowe
\begin{equation} \label{eq:emc2}
	E = mc^2,
\end{equation}
które cytujemy w tekście jako Rów. \eqref{eq:emc2} komendą \texttt{\textbackslash eqref}. Warto zwrócić uwagę, że po otoczeniu \texttt{\textbackslash begin\{equation\} \textbackslash end\{equation\}} \textbf{nie ma nowej linii}. Zapewnia to prawidłowe wyjustowanie tekstu.

Równania wieloliniowe
\begin{align}
	\frac{\pi^2}{6} =&\frac{1}{1^2} + \frac{1}{2^2} + \frac{1}{3^2} + \frac{1}{4^2} + \frac{1}{5^2} + \frac{1}{1^6} + \nonumber\\
	 &+ \frac{1}{7^2} + \frac{1}{8^2} + \frac{1}{9^2} + \frac{1}{10^2} + \ldots 
\end{align}
powinny być odpowiednio wyrównanie, w tym przypadku do pierwsezgo ułamka oraz znaku ,,,$+$''. Jeżeli wszystkie linie opisują 1 równanie, starczy mu przydzielić 1 numer jak w przykładzie powyżej. Jeżeli równań jest kilka, można im nadać niezależne numery.

\section{Definicje, lematy, twierdzenia, przykłady i wnioski}
Definicje, lematy, twierdzenia, przykłady i wnioski piszemy w pracy tak:
\begin{definition}[Martyngał]
Tu piszemy treść definicji martyngału.
\end{definition}
\begin{lemma}[]% w nawiasie kwadratowym można napisać jego nazwę
Tu piszemy treść lematu.
\end{lemma}


\section{Przykładowe cytowania}

Do cytowania używamy komendy \texttt{\textbackslash cite}. W nawiasie klamrowym podajemy klucz, którego użyliśmy w pliku \emph{bibliografia.bib}. Przykład: \cite{einstein} lub \cite[chap. 2]{latexcompanion}.


\section{Przykładowe linki zewnętrzne}

Przykładowy link zewnętrzny wygląda następująco \url{www.google.com}. Można w ten sposób np. linkować do strony GitHuba pracy.

\chapter{Podsumowanie}
Rozdział podsumowanie/dyskusja stanowi konkluzję pracy. Przypominamy co osiągnęliśmy i co z tego wynika. Nie jest on obowiązkowy, ale w większości przypadków lepiej go zamieścić.

\appendix
\chapter{Dodatkowe informacje}
Tu opcjonalnie mogą znajdować się informacje, które stanowią uzupełnienie pracy, np. dodatkowe dane, kod używanych programów, mniej znaczące fragmenty wyprowadzeń, itp.

%%%%%%%%%%%%%%%%%%%%%%%%%%%%%%%%%%%%%%%%%%%%%%%%%%%%%%%%%
% BIBLIOGRAFIA
% W tworzeniu bibliografii najlepiej korzystać z BibTexa, 
% który jest częścią systemu Tex. W naszym przypadku funkcję 
% przechowalni literatury, do której się odwołujemy, pełni 
% plik bibliografia.bib. Nie musimy ręcznie dodawać nowych 
% pozycji do bibliografii. Możemy wejść np. na stronę 
% https://www.doi2bib.org/
% wpisać numer doi odpowiedniej publikacji, skopiować uzyskany 
% tekst, wkleić do pliku bibliografia.bib i skompilować. 
% Gotowe informacje do pliku bibliografia.bib można znaleźć 
% także na https://arxiv.org oraz stronach wydawnictw naukowych.

%%%%%%%%%%%%%%%%%%%%%%%%%%%%%%%%%%%%%%%%%%%%%%%%%%%%%%%%%
\newpage

% styl bibliografii, tutaj unsrt
\bibliographystyle{unsrt}
% w nawiasie klamrowym wpisujemy nazwę pliku z bibliografią w formacie .bib
\bibliografia{bibliografia} 
\end{document}